%-------------------------------------------------------------------------------
%	SECTION TITLE
%-------------------------------------------------------------------------------
\cvsection{Sunumlar \& Yayınlar}
\vspace{-2mm}


%-------------------------------------------------------------------------------
%	CONTENT
%-------------------------------------------------------------------------------
\begin{cventries}

%---------------------------------------------------------
  \cventry
    {YAZAR | ARAŞTIRMACI} % Role
    {\underline{\href{https://www.eucongress.org/_files/ugd/8f4327_0625909f99ad4ff7ac587a9a1fd87003.pdf}{13\textsuperscript{th} International European Conference on Interdisciplinary Scientific Research}}} % Event
    {Tiran, Arnavutluk} % Location
    {May 2026} % Date(s)
    {
      \begin{cvitems} % Description(s)
        \item {Akü üretim hatlarında elektrik ve basınçlı hava tüketimini tahminleyen ML regresyon modelleri (XGBoost, LightGBM, GBM) geliştirdi; kısa vadeli tahminlerde \textbf{R\textsuperscript{2} değerini 0.944'e kadar} çıkardı.}
      \end{cvitems}
    }

%---------------------------------------------------------
  \cventry
    {YAZAR | ARAŞTIRMACI} % Role
{\underline{\href{https://argeinv.mcbu.edu.tr/kongrelerimiz/altinci/ozet/ozet.pdf}{R\&D \& Innovation 2024}}} % Event
    {Manisa, Türkiye} % Location
    {Ara 2024} % Date(s)
    {
      \begin{cvitems} % Description(s)
        \item {Endüstri 4.0 ve IIoT için Veri Entegrasyonu: \textbf{MQTT}, \textbf{OPC UA} ve \textbf{Node.js} ile Unified Namespace Tabanlı Dijital Dönüşüm.}
      \end{cvitems}
    }

%---------------------------------------------------------
  \cventry
    {YAZAR | ARAŞTIRMACI} % Role
{\underline{\href{https://www.ichoracongress.com/hora2021/HORA_2021_abstracts_book.pdf}{Human-Computer Interaction Optimization and Robotic Applications}}} % Event
    {Türkiye} % Location
    {Haz 2021} % Date(s)
    {
      \begin{cvitems} % Description(s)
        \item {Poincaré Grafiği ölçümlerinden \textbf{Makine Öğrenmesi} (topluluk öğrenmesi) yöntemleriyle proses kontrol sistemlerinde hata tespiti ve teşhisi; Tennessee Eastman Sürecinde sınıflandırma doğruluğunu \textbf{\%89,5}'e çıkardı.}
      \end{cvitems}
    }

%---------------------------------------------------------
  \cventry
    {YAZAR | ARAŞTIRMACI} % Role
{\underline{\href{https://www.researchgate.net/profile/Yalcin-Isler-2/publication/357078829_Heart_Sounds_Analysis_and_Classification_Based_on_Long-Short_Term_Memory/links/61c338cdabcb1b520ad8e586/Heart-Sounds-Analysis-and-Classification-Based-on-Long-Short-Term-Memory.pdf}{International Medical Device Conference}}} % Event
    {Antalya, Türkiye} % Location
    {Eyl 2020} % Date(s)
    {
      \begin{cvitems} % Description(s)
        \item {\textbf{Uzun-Kısa Vadeli Bellek (LSTM)} Tabanlı Kalp Sesi Analizi ve Sınıflandırması: fonokardiyogram kayıtlarını Normal, Hırıltılı, Ekstrasistol ve Yapay kategorilerine sınıflandıran bir LSTM modeli geliştirdi; \textbf{\%79,0} doğruluğa ulaştı (Akıllı Sistemler ve Uygulamaları Dergisi, Cilt 3(1), 2020).}
      \end{cvitems}
    }

%---------------------------------------------------------
\end{cventries}
